% Part: first-order-logic
% Chapter: model-theory
% Section: partial-iso

\documentclass[../../../include/open-logic-section]{subfiles}

\begin{document}

\olfileid[pt]{mod}{bas}{pis}
\section{Isomorfismos parciais}

\begin{defn}
  Dadas duas !!{structure}s $\Struct{M}$ e $\Struct{N}$, um
  \emph{isomorfismo parcial} de $\Struct{M}$ em $\Struct{N}$ é uma
  função parcial finita $p$ que recebe argumentos em $\Domain M$ e retorna
  valores em $\Domain N$, e que satisfaz em seu domínio as condições de
  isomorfismo de \olref[iso]{defn:isomorphism}:
  \begin{enumerate}
  \item $p$ é !!{injective};
  \item para toda !!{constant}~$c$: se $p(\Assign{c}{M})$ estiver definida,
    então $p(\Assign{c}{M}) = \Assign{c}{N}$;
  \item para todo !!{predicate} de $n$ lugares $P$: se $a_1$, \dots, $a_n$
    estiverem no domínio de $p$, então $\langle a_1, \dots, a_n\rangle \in
    \Assign P M$ se, e somente se, $\langle p(a_1), \dots, p(a_n) \rangle
    \in \Assign P N$;
  \item para toda !!{function} de $n$ lugares $f$: se $a_1$, \dots, $a_n$
    estiverem no domínio de $p$, então $p(\Assign f M (a_1, \dots,a_n))
    = \Assign f N (p(a_1), \dots, p(a_n))$.
  \end{enumerate}
  Dizer que $p$ é finita significa que $\dom{p}$ é finito.
\end{defn}

Observe que a função vazia~$\emptyset$ é sempre um isomorfismo parcial
entre quaisquer duas !!{structure}s.

\begin{defn}\ollabel{defn:partialisom}
  Duas !!{structure}s $\Struct{M}$ e $\Struct{N}$ são
  \emph{parcialmente isomorfas}, o que se escreve $\Struct{M} \iso[p]
  \Struct{N}$, se, e somente se, houver um conjunto não vazio $I$
  de isomorfismos parciais entre $\Struct{M}$ e $\Struct{N}$
  que satisfaça a propriedade de \emph{ida e volta}:
  \begin{enumerate}
  \item (\emph{Ida}) Para todo $p \in I$ e todo $a \in \Domain M$,
    existe $q \in I$ tal que $p \subseteq q$ e $a$ está
    no domínio de $q$;
  \item (\emph{Volta}) Para todo $p \in I$ e todo $b \in \Domain N$,
    existe $q \in I$ tal que $p \subseteq q$ e $b$ está
    na imagem de $q$.
  \end{enumerate}
\end{defn}

\begin{thm}\ollabel{thm:p-isom1}
  Se $\Struct{M} \iso[p] \Struct{N}$ e $\Struct{M}$ e
  $\Struct{N}$ forem !!{enumerable}, então $\Struct{M} \iso
  \Struct{N}$.
\end{thm}

\begin{proof}
  Como $\Struct{M}$ e $\Struct{N}$ são !!{enumerable}, sejam $\Domain{M} =
  \{a_0, a_1, \ldots \}$ e $\Domain{N} = \{b_0, b_1, \ldots \}$. Partindo
  de um $p_0 \in I$ arbitrário, definimos uma sequência crescente de
  isomorfismos parciais $p_0 \subseteq p_1 \subseteq p_2
  \subseteq \cdots$ do seguinte modo:
  \begin{enumerate}
  \item se $n+1$ for ímpar, digamos $n = 2r$, então, usando a propriedade
    de ida, encontre um $p_{n+1} \in I$ tal que $p_n \subseteq p_{n+1}$
    e $a_r$ esteja no domínio de $p_{n+1}$;
  \item se $n+1$ for par, digamos $n+1 =2r$, então, usando a propriedade
    de volta, encontre um $p_{n+1} \in I$ tal que $p_n \subseteq p_{n+1}$
    e $b_r$ esteja na imagem de $p_{n+1}$.
  \end{enumerate}
Se agora fizermos:
\[
p = \bigcup_{n\ge 0} p_n,
\]
teremos que $p$ é um um isomorfismo entre $\Struct{M}$ e
$\Struct{N}$.
\end{proof}

\begin{prob}
  Mostre detalhadamente que $p$, tal como definida em
  \olref[mod][bas][pis]{thm:p-isom1}, é de fato um isomorfismo.
\end{prob}

\begin{thm}\ollabel{thm:p-isom2}
  Suponha que $\Struct{M}$ e $\Struct{N}$ sejam !!{structure}s para uma
  !!{language} puramente relacional (!!a{language} que contém somente
  !!{predicate}s, e sem !!{function}s nem constantes). Então, se
  $\Struct{M} \iso[p] \Struct{N}$, também $\Struct{M} \elemequiv
  \Struct{N}$.
\end{thm}

\begin{proof}
  Por indução nas !!{formula}s, mostra-se que, se $a_1$, \dots, $a_n$ e
  $b_1$, \dots, $b_n$ forem tais que haja um isomorfismo parcial $p$
  que leve cada $a_i$ a $b_i$, e $s_1(x_i) =a_i$ e $s_2(x_i) =b_i$
  (para $i =1$, \dots,~$n$), então $\Sat{M}{!A}[s_1]$ se,
  e somente se, $\Sat{N}{!A}[s_2]$. O caso $n=0$
  fornece $\Struct{M} \elemequiv \Struct{N}$.
\end{proof}

\begin{rem}
Se houver !!{function}s, o resultado anterior ainda será verdadeiro, mas
será preciso considerar o isomorfismo induzido por $p$ entre a
sub!!{structure} de $\Struct{M}$ gerada por $a_1$, \dots, $a_n$ e a
sub!!{structure} de $\Struct{N}$ gerada por $b_1$, \dots, $b_n$.
\end{rem}

O resultado anterior pode ser ``decomposto'' em etapas estabelecendo-se uma
relação entre o número de quantificadores aninhados em !!a{formula} e
quantas vezes os isomorfismos parciais pertinentes podem ser estendidos.

\begin{defn}
  Para toda !!{formula}~$!A$, o \emph{posto de quantificadores} de $!A$,
  denotado por $\QuantRank{!A} \in \Nat$, é definido recursivamente como
  o maior número de quantificadores aninhados em $!A$. Duas
  !!{structure}s $\Struct{M}$ e $\Struct{N}$ são \emph{$n$-equivalentes},
  o que se escreve $\Struct{M} \elemequiv[n] \Struct{N}$, se coincidirem
  em todas as sentenças de posto de quantificadores menor ou igual a~$n$.
\end{defn}

\begin{prop}\ollabel{prop:qr-finite}
  Seja $\Lang{L}$ uma !!{language} puramente relacional finita, isto é, uma
  !!{language} que contém um número finito de !!{predicate}s e !!{constant}s,
  e sem !!{function}s. Então, para cada $n \in \Nat$, há apenas um
  número finito de !!{sentence}s de primeira ordem na !!{language}
  $\Lang{L}$ que têm posto de quantificadores não maior que $n$, a menos
  de equivalência lógica.
\end{prop}

\begin{proof}
  Por indução em $n$.
\end{proof}

\begin{defn}
  Dada !!a{structure}~$\Struct{M}$, seja $\Domain M^{<\omega}$ o conjunto de
  todas as sequências finitas sobre $\Domain{M}$. Usamos $\mathbf{a},
  \mathbf{b}, \mathbf{c}, \ldots$ para percorrer sequências finitas de
  elementos. Se $\mathbf{a} \in \Domain{M}^{<\omega}$ e $a \in \Domain{M}$, então
  $\mathbf{a}a$ representa a \emph{concatenação} de $\mathbf{a}$ com $a$.
\end{defn}

\begin{defn}
  Dadas !!{structure}s $\Struct{M}$ e $\Struct{N}$, definimos
  relações $I_n \subseteq \Domain M^{<\omega} \times \Domain N^{<\omega}$ entre
  sequências de mesmo comprimento, por recursão em $n$, do seguinte modo:
   \begin{enumerate}
   \item $I_0(\mathbf{a},\mathbf{b})$ se, e somente se, $\mathbf{a}$ e
     $\mathbf{b}$ satisfizerem as mesmas !!{formula}s atômicas em $\Struct{M}$
     e $\Struct{N}$; isto é, se $s_1(x_i) = a_i$ e $s_2(x_i) =
     b_i$, e $!A$ for atômica com todas as !!{variable}s entre
     $x_1$, \dots,~$x_n$, então $\Sat{M}{!A}[s_1]$ se,
     e somente se,~$\Sat{N}{!A}[s_2]$.
   \item $I_{n+1} (\mathbf{a},\mathbf{b})$ se, e somente se, para todo
     $a\in \Domain M$ houver um $b\in \Domain N$ tal que $I_n
     (\mathbf{a}a,\mathbf{b}b)$, e vice-versa.
   \end{enumerate}
\end{defn}


\begin{defn}
  Escrevemos $\Struct{M} \approx_n \Struct{N}$ se
  $I_n(\emptyseq,\emptyseq)$ valer para $\Struct{M}$ e
  $\Struct{N}$ (onde $\emptyseq$ é a sequência vazia).
\end{defn}

\begin{thm}\ollabel{thm:b-n-f}
  Seja $\Lang{L}$ uma !!{language} puramente relacional. Então $I_n
  (\mathbf{a},\mathbf{b})$ implica que, para todo $!A$ tal que
  $\QuantRank{!A} \le n$, temos $\Sat{M}{!A}[\mathbf{a}]$ se, e
  somente se, $\Sat{N}{!A}[\mathbf{b}]$ (onde, novamente, $\mathbf{a}$
  satisfaz $!A$ se qualquer $s$ tal que $s(x_i) = a_i$ satisfizer
  $!A$). Além disso, se $\Lang{L}$ for finita, a recíproca também vale.
\end{thm}

\begin{proof}
  A demonstração de que $I_n(\mathbf{a},\mathbf{b})$ implica que
  $\mathbf{a}$ e $\mathbf{b}$ satisfazem as mesmas !!{formula}s de
  posto de quantificadores não maior que $n$ é feita por uma indução
  simples em $!A$. Para a recíproca, procedemos por indução em $n$, usando
  \olref{prop:qr-finite}, que garante que, para cada $n$,
  há no máximo finitamente muitas !!{formula}s não equivalentes desse
  posto de quantificadores.

  Para $n=0$, a hipótese de que $\mathbf{a}$ e $\mathbf{b}$ satisfazem
  as mesmas !!{formula}s sem quantificadores implica que elas satisfazem
  as mesmas atômicas, de modo que $I_0(\mathbf{a},\mathbf{b})$.

  Para o caso $n+1$, suponha que $\mathbf{a}$ e $\mathbf{b}$
  satisfaçam as mesmas !!{formula}s de posto de quantificadores não maior
  que $n+1$; para mostrar que $I_{n+1}(\mathbf{a},\mathbf{b})$, basta
  mostrar que, para cada $a \in \Domain M$, há um $b \in
  \Domain N$ tal que $I_n(\mathbf{a}a,\mathbf{b}b)$ e, pela hipótese
  de indução, basta novamente mostrar que, para cada $a \in
  \Domain M$, há um $b \in \Domain N$ tal que $\mathbf{a}a$ e
  $\mathbf{b}b$ satisfazem as mesmas !!{formula}s de posto de
  quantificadores não maior que $n$.

  Dado $a \in \Domain M$, seja $!T^a_n$ conjunto de !!{formula}s
  $!B(x,\mathbf{y})$ de posto não maior que $n$ satisfeitas por
  $\mathbf{a}a$ em $\Struct{M}$; $!T^a_n$ é finito, portanto podemos
  supor que seja uma única !!{formula} de primeira ordem. Segue-se que
  $\mathbf{a}$ satisfaz $\lexists[x][!T^a_n(x,\mathbf{y})]$, que
  tem posto de quantificadores não maior que $n+1$. Por hipótese,
  $\mathbf{b}$ satisfaz a mesma !!{formula} em $\Struct{N}$, de modo que
  há um $b \in \Domain N$ tal que $\mathbf{b}b$ satisfaz
  $!T^a_n$; em particular, $\mathbf{b}b$ satisfaz as mesmas
  !!{formula}s de posto de quantificadores não maior que $n$ que
  $\mathbf{a}a$. De modo semelhante, mostra-se que, para todo $b \in \Domain N$,
  há $a\in \Domain M$ tal que $\mathbf{a}a$ e $\mathbf{b}b$
  satisfazem as mesmas !!{formula}s de posto de quantificadores não maior
  que $n$, o que conclui a demonstração.
\end{proof}

\begin{cor}\ollabel{cor:b-n-f}
  Se $\Struct{M}$ e $\Struct{N}$ forem !!{structure}s puramente relacionais
  em uma !!{language} finita, então $\Struct{M} \approx_n\Struct{N}$ se,
  e somente se, $\Struct{M} \elemequiv[n] \Struct{N}$. Em particular,
  $\Struct{M} \elemequiv \Struct{N}$ se, e somente se, para cada $n$,
  $\Struct{M} \approx_n \Struct{N}$ .
\end{cor}

\end{document}
